\documentclass{article}

\usepackage[main, preprint]{neurips_2026}
\usepackage[utf8]{inputenc} % allow utf-8 input
\usepackage[T1]{fontenc}    % use 8-bit T1 fonts
\usepackage{hyperref}       % hyperlinks
\usepackage{url}            % simple URL typesetting
\usepackage{booktabs}       % professional-quality tables
\usepackage{amsfonts}       % blackboard math symbols
\usepackage{nicefrac}       % compact symbols for 1/2, etc.
\usepackage{microtype}      % microtypography
\usepackage{xcolor}         % colors
\usepackage{graphicx}
\usepackage{float}
\usepackage{amsmath}

\title{Greenbyte}


\author{%
  Alexis Marez \\
  Department student belongs to\\
  University of Colorado Boulder\\
  \texttt{alma2309@colorado.edu} \\
  % examples of more authors
  \And
  James Vallery\\
  Department student belongs to\\
  University of Colorado Boulder\\
  \texttt{java2626@colorado.edu} \\
  % \AND
  % Coauthor \\
  % Affiliation \\
  % Address \\
  % \texttt{email} \\
  % \And
  % Coauthor \\
  % Affiliation \\
  % Address \\
  % \texttt{email} \\
  % \And
  % Coauthor \\
  % Affiliation \\
  % Address \\
  % \texttt{email} \\
}


\begin{document}


\maketitle


\begin{abstract}
  Managing microclimates in small-scale greenhouses is traditionally hindered by the high cost of sensor arrays and the long duration required for manual data collection. We present a methodology for creating a personalized greenhouse climate predictor using a "digital twin" approach. By calibrating the GreenLight simulator via Nelder-Mead optimization to match the physical thermal envelope of a specific residential greenhouse, we generated dataset spanning 500 diverse environmental scenarios. This data was used to train a Neural Network capable of predicting internal temperature, relative humidity, and operating costs based on external weather inputs. Our results demonstrate that simulation-to-reality pipelines can bypass data collection roadblocks, providing smaller growers with accurate, personalized predictive tools that achieve a target MAE of <2.0 F.
\end{abstract}


\section{Introduction}

Managing a greenhouse requires control over different variables such as temperature, relative humidity, and $CO_2$ levels to ensure optimal plant growth and cost. However, many small-scale greenhouse owners lack the expensive sensors and acuators required for advanced automation. The problem we address is creating an accurate predictive model that is tailored to a specific greenhouse without requiring months of manual data collection. By using the GreenLight simulator, we can bypass initial data collection hurdles by simulating the specific physical constraints of the student’s greenhouse. This project is important because it demonstrates how digital twins and simulation can lower the barrier to entry for smart-agriculture technology. 

\subsection{Related Work}
Current research in greenhouse climate modeling often relies on either purely physical models (like the mechanics used in GreenLight) or data-driven models trained on large-scale industrial datasets. Our approach builds on these by bridging the gap: using the physical accuracy of GreenLight to simulate a real world greenhouse, then train a personalized machine learning model that adapts to the unique variables of the settup.

\section{Learning and Project Goals}

Our primary objective is to build a tool that can predict variables such as internal temperature, relative humidity, VPD (vapor pressure deficit), and cost with high accuracy based on external weather inputs and current greenhouse conditions.

Learning Goals: We aim to understand the pipeline of generating synthetic data from simulators and the nuances of training time-series models for environmental forecasting. 

Quantitative Objectives: We hope to achieve a Root Mean Squared Error (RMSE) 2.0°C (3.6°F) for temperature predictions and a test error of at most 0.15 normalized RMSE across all metrics (temperature, relative humidity, and heating energy consumption).

\section{Background}

\subsection{Theory Gaussian Processes}

Write about a theoretical aspect of your project that you chose to understand more deeply.
For example, you could choose to understand cross entropy loss, ridge regression, Gaussian processes, or any fundamental component that warrants deeper understanding for your project.
Explain the concept clearly, with formal mathematical notation and any relevant results, with citations.
For example, for Gaussian processes, you could state the model formally and cite conditions on the kernel for the model to be continuous with probability 1.  For cross entropy, you could prove that it is "proper", i.e.\ the best model does minimal the loss in expectation, or useful properties related to Shannon information theory.


\section{Method}
The following section details our experimental pipeline and model architecture. We describe the transition from physical simulation to a data-driven predictive model, justifying the design choices made to capture the greenhouse's unique thermodynamics.

\begin{itemize}
    \item \textbf{Simulation Pipeline:} We utilized the GreenLight simulator to model the specific physical parameters of one of the researchers personal greenhouse. This effectively created a digital twin of the researcher greenhouse. We did this by calibrating the greenhouses thermal envelope. The values were tuned via Nelder-Mead optimization against the Verdify dataset (real world greenhouse data) \begin{table}[h]                       
  \centering                                                                                                                      
      \label{tab:greenlight-params}               
      \begin{tabular}{llll}                                                                                                              
      \hline                                                                                                                             
      \textbf{Parameter} & \textbf{Symbol} & \textbf{Value} & \textbf{Notes} \\                                                          
      \hline                                                                                                                             
      \multicolumn{4}{l}{\textit{Physical structure}} \\                                                                                 
      Cover area          & $A_{\text{cov}}$       & $74.8\ \text{m}^2$          & North wall excluded (attached to house) \\
      Floor area          & $A_{\text{flr}}$       & $34.1\ \text{m}^2$          & \\                                                    
      \hline                                      
      \multicolumn{4}{l}{\textit{Mechanical systems}} \\                                                                                 
      Forced ventilation  & $\phi_{\text{vent}}$   & $2.31\ \text{m}^3/\text{s}$ & 2 $\times$ 2450 CFM exhaust fans \\                   
      LED lamp intensity  & $\theta_{\text{lamp}}$ & $42.4\ \text{W/m}^2$        & 49 $\times$ Barrina T8, 1446 W total \\               
      \hline                                                        
    
    
      \multicolumn{4}{l}{\textit{Calibrated thermal envelope}} \\
        Infiltration leakage & $c_{\text{leakage}}$ & $1.50 \times 10^{-3}$ & Tuned; default $3.0 \times 10^{-4}$ \\
        Cover conductivity   & $\lambda_{\text{Rf}}$ & $0.057$ W/(m$\cdot$K) & Tuned; default $0.021$ W/(m$\cdot$K) \\
        \hline                                                                                                                           
      \multicolumn{4}{l}{\textit{External weather}} \\                                                                                    
      Station & \multicolumn{3}{l}{Boulder Muni AP, WMO 720533, elev.\ 1612 m (Boulder, CO TMYx)} \\                                     
      \hline                                                                                                                             
      \end{tabular}                                                                                                                      
      \end{table}   
    \item \textbf{Data Acquisition:} The simulator was run for 30 days per sample across 500 random sampled parameter combinations, with simulation start days distributed uniformly across the calendar year so all seasons are represented. This generated a CSV dataset of 500 rows, each representing one 30-day simulation run, with 20 input/output columns. Key output variables measured include mean air temperature ($\bar{T}{\text{air}}$, °C), minimum air temperature, mean relative humidity, mean canopy temperature, heating energy ($\text{MJ/m}^2$), lighting energy ($\text{MJ/m}^2$), water consumption (L/m²), and total operating cost (USD/m²).
    \item \textbf{Model Architecture:} We implemented a Neural Network to map these simulated inputs to the desired output variables. Our csv data set was split into training, validation, and test datasets with the percentage being 80/20/20 respectively. We had it run for 100 epochs \textbf{(still working on finalzing the model by experimenting with these variables)} while changing learning rate, batch size, number of hidden layers, and activation functions until we measured the model to converge \textbf{(again we plan to be more specific once we find the optimal hyperparameters)}. 
    \item \textbf{Design Justification:} This architecture was chosen for a number of reasons. To start, GreenLight provided us with a way to create a large greenhouse dataset with many variables, both output and input. When we compare the number of variables in each entry of our dataset to the variables in public datasets like Kaggle, we have hundreds more to predict/edit to ensure we have an accurate model specifically tailored to our greenhouse. 

    The Neural Network was chosen for its ability to handle the complexity our dataset has to offer. With this we were able to feature engineer input variables to create more accurate predictions. 
    
\end{itemize}

\subsection{Technical Formalism}

\section{Discussion}
 So far, the most important part of the calibration hasn't been a number, its being a diagnosis of how the greenhouse behaves in different external environments. Early calibration runs produced a simulation where the heater never fired despite outdoor temperatures in the teens F and a heater runtime of 67.3 percent. 
 
 A temperature MAE of 5.73°F looked like a thermal envelope problem, it was not. GreenLight's default heating setpoint of 58°F sat below the real greenhouse's minimum indoor temperature of 62.4°F, so the simulated thermostat never called for heat. The optimizer, given no other explanation for the cold bias, began warping the infiltration coefficient $c_{\text{leakage}}$ upward to compensate, this produced a parameter value $5\times$ the physical default. 
 
 This would have gone undetected if we had only reported temperature MAE. Including a heater runtime fraction as a scored metric made the failure immediately obvious and shifted the optimizers approach.

Some questions to consider then:
\begin{enumerate}

    \item \textbf{How do we know that the simulation is actually reflective of the physical world, rather than GreenLight simply tuning an unknown parameter?}

    This is fundamentally a question of \emph{model validity rather than model performance}. A simulation such as GreenLight can be calibrated to match observed outputs (e.g., energy use), but this does not guarantee that the internal parameters correspond to real-world physical processes.

    In other words, an optimizer can ``turn knobs'' to reduce error without those knobs representing true causal relationships. This creates a risk of overfitting to observed outcomes rather than capturing real values.

    To address this, simulation must rely on more than a single error metric. It requires:
    \begin{itemize}
        \item Comparing multiple outputs (not just one target)
        \item Validating against independent datasets
        \item Ensuring parameter values remain physically plausible
    \end{itemize}

    Agreement with real-world data is therefore not sufficient evidence of correctness.

    \item \textbf{How do we measure simulation effectiveness or calibration, given that numerical accuracy alone is insufficient?}

    Simulation effectiveness cannot be captured by a single numerical metric because greenhouse systems are inherently multi-dimensional. A model may accurately predict VPD while misrepresenting energy consumption or CO\textsubscript{2} dynamics.

    Therefore, effectiveness should be evaluated using a multi-objective framework, including:
    \begin{itemize}
        \item Predictive accuracy across multiple outputs (VPD, energy, CO\textsubscript{2})
        \item Stability under small perturbations of inputs
        \item Physical plausibility of intermediate variables
        \item Generalization to unseen scenarios
    \end{itemize}

    A strong simulation is not only accurate, but also consistent, robust, and interpretable.

    \item \textbf{Did your techniques or approach match your goals?}

    The goal of this project was not only to fit a model, but to explore how simulation data could be used to optimize greenhouse performance. Using GreenLight enabled the generation of large, structured datasets with many controllable variables, which aligned well with this objective.

    However, this introduced a tradeoff: while simulated data provides scalability and control, it raises questions about realism. As a result, the approach supported model experimentation and optimization, but only partially addressed real-world applicability.

    This reflects a broader tension between scalability and realism in simulation-based approaches.

    \item \textbf{Can you think of ways your analysis or tool could fail (e.g., user interaction, system-level issues, or adversarial data)?}

    Potential failure modes exist at multiple levels:

    \begin{itemize}
        \item \textbf{User-level:} Users may interpret outputs as deterministic predictions rather than probabilistic estimates, leading to poor decisions.
        \item \textbf{System-level:} The model may optimize for a single objective (e.g., yield) while neglecting others (e.g., energy efficiency).
        \item \textbf{Data-level:} Reliance on simulated data may introduce bias or unrealistic assumptions. Additionally, adversarial or extreme inputs may produce physically implausible outputs.
    \end{itemize}

    These issues highlight that the system operates within a broader decision-making context where errors can propagate.

    \item \textbf{What roadblocks did you encounter, and how did you adapt your approach?}

    One major challenge was the lack of sufficiently rich real-world greenhouse datasets. Available datasets were often too small or lacked the necessary feature complexity for meaningful modeling.

    This led to a pivot toward simulation using GreenLight, which provided greater scale and control. However, this introduced new challenges related to realism and validation.

    As a result, the project evolved from a purely data-driven modeling approach into an exploration of simulation-based learning and its limitations, shifting the focus from prediction accuracy to questions of validity and interpretability.

\end{enumerate}


\subsection{Response to Feedback [required]}

\begin{itemize}
    \item \textbf{Exploratory Data Analysis and Feature Engineering:} We were suggested to perform Canonical Correlation Analysis (CCA) to find relationships we variables. We took the brute force path of experimenting with different input variables and simply crunching numbers using simple regression models to find relationships between variables.
    
    \item \textbf{Model Progression:} We were also told to "start with simple models" which led us to initially implement a Linear Regression baseline and an XGBoost model. While the linear model provided a benchmark of model performance, the XGBoost model showed us the non-linear nature of the greenhouse. This progression helped us transition to a more complex Neural Network, because simpler models struggled to capture the dependencies of our output variables with the many complex input variables.
    
    \item \textbf{Preprocessing and Dataset Scaling:} Taking the advice that "more data is better," we increased our GreenLight simulation runs to ensure a diverse range of seasonal weather patterns. This allowed us to perform more robust preprocessing, including handling outliers in the synthetic sensor data and normalizing inputs for the neural network.
\end{itemize}

\subsection{Ethics \& Impact [required]}
As machine learning becomes increasingly integrated into agricultural systems, we must consider the broader implications of these models. Our greenhouse predictor presents specific ethical considerations regarding data privacy, model reliability, and accessibility.

\begin{itemize}
    \item \textbf{Privacy and Data Collection:} The data used in this project was collected from a personal greenhouse belonging to a researcher. Data such as weather is higlhy available and does pose a signficant risk to privacy. However data about detailed heaters and lighting schedules can  reveal patterns of human activity or residency. If this was a real application to be put out to production or to be made as a public repository, we would ensure to not share geographical data and we would also not require personal information to get started on the use of our project.
    
    \item \textbf{Risk of Harmful Outputs:} A primary ethical concern is the risk of "false confidence" in model predictions. If again this project would be deployed, then a grower relying solely on our model to manage a physical greenhouse could suffer great equipment, crop, and total cost losses if an inaccurate prediction were to occur. This is particularly dangerous in extreme weather conditions where the model may not have sufficient training data to predict a freeze event or a dry event.
    
    \item \textbf{Stakeholder Impact:} This technology empowers small-scale hobbyists and researchers by lowering the cost of "smart" climate control. However, we must ensure that such tools do not create a "digital divide" where only those with high-performance computing resources can access personalized optimization.
\end{itemize}

\subsection{Mitigation and Responsibility}
To address potential harms, we would advocate for using this model as a decision-support tool rather than an autonomous controller. This way the risk of biased or incorrect model outputs leading to physical damage is significantly reduced, since users would be encouraged to only leverage the models output, not relying on it. 


\section{Conclusion}
This project successfully demonstrated a pipeline for personalized agricultural forecasting by bridging physical simulation and machine learning. By creating a digital twin of a specific residential greenhouse and calibrating its thermal envelope, we generated a robust synthetic dataset that allowed us to train a predictive Neural Network without months of manual sensing.

Our findings highlight that while physical simulators are essential for data generation, the model's predictive power relies heavily on correct feature engineering and an understanding of the simulator's underlying assumptions. This approach matters because it successfully eases access to smart-farming technology, allowing individual growers to optimize their energy usage and crop health through personalized models.

\subsection*{AI Disclosure [required]}

As for the use of AI we have used it in aiding the development of both the data and the predictive model we used. More specifically AI was used to help understand the GreenLight repository so we can then properly run it with the desired variables. We used Gemini to give suggestions on which models would be the best fit to our data set. We used AI to help parameterize the greenhouse into variables and methods that green light can understand. Additionally it wrote and optimized calibration scripts to sweep over Verdiy csv's files to calculate MAE and RMSE using Nelder-Mead optimizer.



\subsection{Citations}
\bibliography{references}                              \textbf{Section that we need to finish in the final draft.}                                                                                                            \bibliographystyle{plainnat}    
                                                                                                                                                                                                                                                                           
\end{document}